\chapter{Kết luận và Hướng phát triển}
\label{ch:ket-luan}

\section{Những kết quả đạt được của đề tài}
Đề tài "Hệ thống hỗ trợ sinh viên - phân hệ quản lý tài chính" đã hoàn thành đầy đủ các mục tiêu đề ra ban đầu cả về nghiên cứu lý thuyết lẫn ứng dụng thực tiễn. Mức độ tin cậy và phương pháp kiểm chứng cho từng thành quả cụ thể được xác định như sau:
\begin{itemize}
    \item \textbf{Về mặt kỹ thuật phần mềm (Độ tin cậy cao):} Hiện thực hóa thành công một kiến trúc Client-Server phân tán nhiều tầng hiện đại, tích hợp trơn tru giữa ứng dụng di động React Native, trang quản trị React JS, Backend NestJS và máy chủ AI FastAPI. Kết quả này được xác minh thông qua các đợt kiểm thử tích hợp hệ thống (Integration Testing) và chạy thử nghiệm thực tế trên các thiết bị di động vật lý cũng như trình duyệt web, đảm bảo sự ổn định của cơ chế hàng đợi BullMQ + Redis đối với luồng dữ liệu bất đồng bộ và đồng bộ thời gian thực.
    \item \textbf{Về mặt quản lý tài chính (Độ tin cậy cao):} Số hóa thành công mô hình quản lý tài chính 6 lọ nổi tiếng, phát triển cơ chế tự động phân chia thu nhập, lập lịch giao dịch định kỳ và đảm bảo tính toàn vẹn số dư động thông qua lịch sử giao dịch thực tế. Tính năng này được kiểm chứng qua các kịch bản kiểm thử chức năng (Functional Testing) đối với các giao dịch thu chi thực tế và đối soát số dư tự động, đảm bảo không xảy ra hiện tượng lệch số dư.
    \item \textbf{Về mặt học bổng (Độ tin cậy cao):} Hoàn thiện quy trình học bổng trực tuyến từ công bố chương trình, cấu hình điều kiện và tài liệu yêu cầu, sinh viên lưu nháp hoặc nộp hồ sơ, đến nhà tài trợ xét duyệt và phản hồi kết quả. Kết quả được kiểm chứng thông qua việc kiểm thử các luồng nghiệp vụ chuẩn và các ràng buộc logic, đảm bảo tính nhất quán của trạng thái hồ sơ và ngăn chặn hoàn toàn việc nộp trùng lặp hồ sơ cho cùng một sinh viên và học bổng.
    \item \textbf{Về mặt ứng dụng AI (Độ tin cậy trung bình):} Xây dựng hệ tác nhân AI trợ lý ảo hỗ trợ trả lời câu hỏi tài chính của sinh viên thời gian thực qua stream SSE và huấn luyện bộ phân loại giao dịch tự động từ ngôn ngữ tự nhiên. Độ chính xác trung bình của bộ phân loại đạt \textbf{91.8\%}, đây là số liệu ước tính qua quan sát và đánh giá thủ công trên tập dữ liệu thử nghiệm quy mô nhỏ (100 mẫu câu truy vấn) trong môi trường phát triển; kết quả này cần được tiếp tục đánh giá trên các tập benchmark lớn hơn và chuẩn hóa hơn.
    \item \textbf{Về mặt quản trị dữ liệu học thuật (Độ tin cậy cao):} Hoàn thiện luồng quản lý trường đại học, khoa/chương trình đào tạo và môn học trên Webadmin; hỗ trợ import hàng loạt, lọc dữ liệu và đồng bộ danh mục môn học sang Mobile App để sinh viên sử dụng trong bảng điểm cá nhân. Kết quả được kiểm thử và xác thực thông qua việc kiểm tra cấu trúc dữ liệu cơ sở dữ liệu sau đồng bộ và khả năng hiển thị tương thích trên thiết bị di động.
    \item \textbf{Về mặt bảng điểm sinh viên (Độ tin cậy cao):} Hoàn thiện chức năng cho phép sinh viên theo dõi GPA, tổng tín chỉ, mục tiêu tín chỉ, thêm/sửa/xóa điểm từng môn và import điểm hàng loạt từ file TXT/CSV. Cơ chế tính toán GPA có trọng số theo tín chỉ đã được kiểm chứng chéo với các công thức tính toán GPA chuẩn của các trường đại học lớn tại Việt Nam, đảm bảo độ chính xác của kết quả đầu ra.
\end{itemize}

\subsection{Định vị sản phẩm và phân tích sự khác biệt hóa}
Để làm rõ giá trị cốt lõi và định vị thực tế của hệ thống Student360 trong hệ sinh thái ứng dụng hỗ trợ người dùng, nhóm phát triển tiến hành đối chiếu các nhóm chức năng của hệ thống với ba ứng dụng quản lý tài chính phổ biến hiện nay: Money Lover, Sổ Thu Chi MISA và Spendee. Cần lưu ý rằng, ba ứng dụng được lựa chọn đối chứng là các giải pháp chuyên biệt cho lĩnh vực quản lý tài chính cá nhân dành cho đại chúng, trong khi Student360 được định hình là một giải pháp lai (hybrid) tích hợp giữa tài chính và học thuật dành riêng cho sinh viên. Do đó, việc so sánh được chia thành hai nhóm nhằm phân tích rõ tính tương đồng và sự khác biệt: (1) Tính năng tài chính cốt lõi (so sánh cùng phạm vi để đánh giá mức độ đáp ứng nhu cầu tài chính); (2) Tính năng hỗ trợ học bổng và học thuật (thể hiện sự khác biệt hóa về định vị sản phẩm, vốn không nằm trong mục tiêu phát triển của các đối thủ cạnh tranh).

\textbf{Nhóm 1: Các tính năng Tài chính tương đồng (so sánh cùng phạm vi)}

\begin{table}[H]
    \centering
    \caption{Đối chiếu tính năng Tài chính cốt lõi giữa các giải pháp}
    \label{tab:comparison-finance}
    \begin{tabular}{|L{4.2cm}|C{2.6cm}|C{2.8cm}|C{2.8cm}|C{2.6cm}|}
        \hline
        \textbf{Tính năng Tài chính} & \textbf{Money Lover} & \textbf{Sổ Thu Chi MISA} & \textbf{Spendee} & \textbf{Student360} \\
        \hline
        Tự động phân bổ ví 6 lọ & Thủ công & Thủ công & Thủ công & \textbf{Tự động} \\ \hline
        Nhận diện giao dịch AI bằng NLP & Tiếng Anh & Không áp dụng & Không áp dụng & \textbf{Tiếng Việt} \\ \hline
        Trợ lý ảo tư vấn tài chính & Không & Không & Không & \textbf{Có} \\ \hline
        Lập lịch giao dịch định kỳ & Có & Có & Có & \textbf{Có} \\ \hline
        Hàng đợi xử lý thông báo đẩy & Không cấu hình & Không & Không & \textbf{BullMQ + Redis} \\ \hline
    \end{tabular}
\end{table}

\textit{Phân tích Nhóm 1:} Trong phạm vi quản lý tài chính cá nhân cơ bản, Student360 đáp ứng tốt các nhu cầu cốt lõi tương đương các giải pháp thương mại như Money Lover, MISA, và Spendee. Đồng thời, hệ thống ứng dụng các kỹ thuật như tự động phân bổ thu nhập theo mô hình 6 lọ và xử lý giao dịch tự động bằng NLP tiếng Việt nhằm tối ưu hóa trải nghiệm riêng cho đối tượng người dùng trẻ tuổi.

\textbf{Nhóm 2: Sự khác biệt hóa về Tính năng Học bổng và Học thuật}

\begin{table}[H]
    \centering
    \caption{Phân tích sự khác biệt hóa về tính năng giáo dục (ngoài phạm vi của đối thủ)}
    \label{tab:comparison-education}
    \begin{tabular}{|L{4.2cm}|C{2.6cm}|C{2.8cm}|C{2.8cm}|C{2.6cm}|}
        \hline
        \textbf{Tính năng Giáo dục} & \textbf{Money Lover} & \textbf{Sổ Thu Chi MISA} & \textbf{Spendee} & \textbf{Student360} \\
        \hline
        So khớp học bổng thông minh & Không áp dụng & Không áp dụng & Không áp dụng & \textbf{Có} \\ \hline
        Quy trình học bổng trực tuyến & Không áp dụng & Không áp dụng & Không áp dụng & \textbf{Có} \\ \hline
        Đồng bộ dữ liệu nền đào tạo & Không áp dụng & Không áp dụng & Không áp dụng & \textbf{Có} \\ \hline
        Quản lý dữ liệu học thuật chuẩn hóa & Không áp dụng & Không áp dụng & Không áp dụng & \textbf{Có} \\ \hline
        Bảng điểm cá nhân và GPA theo tín chỉ & Không áp dụng & Không áp dụng & Không áp dụng & \textbf{Có} \\ \hline
    \end{tabular}
\end{table}

\textit{Phân tích Nhóm 2:} Các tính năng thuộc nhóm này không nằm trong phạm vi mục tiêu sản phẩm của Money Lover, MISA hay Spendee. Sự hiện diện của các tính năng này trên Student360 nhằm minh chứng cho hướng đi khác biệt của đề tài: không cố gắng thay thế hoàn toàn các ứng dụng tài chính chuyên nghiệp trên thị trường đại chúng, mà tập trung giải quyết bài toán giao thoa giữa tài chính (học phí, chi tiêu sinh hoạt) và học tập (học bổng, kết quả học tập) của sinh viên.

\textbf{Kết luận định vị sản phẩm:} Việc đặt Student360 cạnh các giải pháp tài chính thương mại cho thấy hệ thống không chỉ đáp ứng tốt các nhu cầu quản lý chi tiêu cơ bản, mà quan trọng hơn là đạt được sự định vị chuyên biệt cao độ cho đối tượng đích thông qua hệ sinh thái giáo dục tích hợp. Sự kết hợp giữa quản lý chi tiêu 6 lọ, công cụ theo dõi học tập cá nhân và cổng thông tin - quy trình học bổng trực tuyến giúp Student360 tạo lập một phân khúc sản phẩm riêng biệt, hỗ trợ toàn diện cuộc sống và học tập của sinh viên.

\section{Hạn chế của hệ thống}
Mặc dù đạt được những kết quả khả quan, hệ thống vẫn tồn tại một số điểm hạn chế cần được cải thiện trong tương lai:
\begin{enumerate}
    \item \textbf{Độ trễ khi gọi LLM:} Thời gian phản hồi trọn vẹn của AI Chatbot phụ thuộc hoàn toàn vào tốc độ sinh token của Gemini API qua mạng internet, đo thực tế trung bình khoảng 4 giây và có thể lên tới hơn 10 giây với các câu hỏi phức tạp cần gọi nhiều tool (xem Bảng~\ref{tab:ai-e2e-results}). Cơ chế lưu snapshot chỉ áp dụng cho phần nhận định tài chính định kỳ (AI Insights), chưa áp dụng được cho hội thoại chatbot vì mỗi câu hỏi của người dùng về bản chất là duy nhất.
    \item \textbf{Giới hạn của cơ chế gợi ý học bổng:} Cơ chế hiện tại hoạt động dựa trên phân tích văn bản tĩnh (Jaccard Token + Bigram + SequenceMatcher), chưa có khả năng nắm bắt ý nghĩa ngữ nghĩa sâu (Semantic Understanding). Kết quả là thông tin trong hồ sơ có từ vựng khác nhau nhưng cùng ý nghĩa (ví dụ: ``kỹ thuật phần mềm'' và ``CNTT'', ``công nghệ thông tin'') có thể không được nhận diện là tương đương trong quá trình gợi ý.
    \item \textbf{Kịch bản kiểm thử dữ liệu:} Dữ liệu học bổng và hồ sơ học thuật sử dụng trong thử nghiệm hiện tại chủ yếu là dữ liệu giả lập mô phỏng. Hệ thống chưa được kiểm thử diện rộng trên quy mô dữ liệu thực tế hàng chục nghìn hồ sơ sinh viên thật từ phòng đào tạo các trường đại học.
    \item \textbf{Sự phụ thuộc vào bên thứ ba:} Cơ chế thông báo di động phụ thuộc lớn vào sự ổn định của dịch vụ đám mây Google FCM và kết nối mạng nền của thiết bị di động.
    \item \textbf{Chưa liên kết tài khoản ngân hàng thực tế:} Luồng ghi nhận thu chi trong hệ thống hiện tại vẫn phụ thuộc vào hành vi nhập liệu thủ công của sinh viên hoặc phê duyệt đề xuất từ trợ lý ảo AI, chưa liên kết trực tiếp với tài khoản ngân hàng hoặc ví điện tử để tự động hóa việc đồng bộ dòng tiền thực tế.
    \item \textbf{Một số thao tác tài chính tổng hợp chưa được bọc trong Database Transaction tường minh:} Cụ thể là luồng phân bổ thu nhập vào 6 hũ (\texttt{distributeIncome}) và luồng xóa hũ kèm giao dịch bù trừ hiện thực hiện các lệnh ghi độc lập thay vì một \texttt{QueryRunner} Transaction chung, nên vẫn còn rủi ro dữ liệu không nhất quán nếu có lỗi xảy ra giữa chừng (xem chi tiết Chương 4).
    \item \textbf{Thiếu bộ benchmark định lượng chính thức cho các thành phần AI và hạ tầng:} Các số liệu về độ chính xác phân loại giao dịch, độ chính xác của cơ chế gợi ý học bổng và khả năng chịu tải của hàng đợi thông báo đẩy FCM hiện dựa trên quan sát và ước tính trong quá trình phát triển, chưa có bộ dữ liệu gán nhãn cố định và script đo lường tự động, có thể tái lập (xem Bảng~\ref{tab:ai-classify-eval}, Bảng~\ref{tab:ai-matching-eval} và Bảng~\ref{tab:fcm-latency-eval}).
\end{enumerate}

\section{Chiến lược phát triển Now--Next--Later}

Để khắc phục các hạn chế hiện tại và nâng cao khả năng ứng dụng của hệ thống trong thực tế, nhóm đề xuất lộ trình phát triển theo ba giai đoạn. Thứ tự ưu tiên được xây dựng dựa trên nguyên tắc hoàn thiện độ tin cậy của hệ thống trước khi mở rộng các tính năng tự động hóa và nâng cao năng lực AI.

\subsection{Now: Nâng cao độ tin cậy hệ thống}

\begin{enumerate}
    \item \textbf{Bổ sung Database Transaction cho các thao tác tài chính tổng hợp:}
    \begin{itemize}
        \item \textbf{Độ ưu tiên:} Cao.
        \item \textbf{Nguồn lực dự kiến:} 01 Backend Developer, môi trường kiểm thử Staging.
        \item \textbf{Sự phụ thuộc:} Không có.
        \item \textbf{Rủi ro và Giải pháp:} Rủi ro xảy ra hiện tượng khóa chết (deadlock) khi đồng thời ghi dữ liệu ví tài chính. Giải pháp là áp dụng cơ chế khóa bi quan (pessimistic locking) và sắp xếp thứ tự truy vấn các bảng liên quan một cách nhất quán trong transaction.
        \item \textbf{KPI và Tiêu chí hoàn thành:} 100\% các luồng tài chính phức tạp (như phân bổ ví 6 lọ, giao dịch bù trừ khi xóa ví) được thực thi trong Transaction; không xảy ra hiện tượng lệch số dư khi giả lập lỗi mất kết nối cơ sở dữ liệu giữa chừng.
    \end{itemize}

    \item \textbf{Xây dựng bộ benchmark định lượng cho các thành phần AI:}
    \begin{itemize}
        \item \textbf{Độ ưu tiên:} Cao.
        \item \textbf{Nguồn lực dự kiến:} 01 AI/ML Engineer, bộ dữ liệu khoảng 500 mẫu hội thoại/phân loại giao dịch được gán nhãn thủ công từ dữ liệu sinh viên thực tế.
        \item \textbf{Sự phụ thuộc:} Không có.
        \item \textbf{Rủi ro và Giải pháp:} Dữ liệu gán nhãn có thể thiếu khách quan hoặc có sai sót. Giải pháp là thực hiện quy trình kiểm tra chéo giữa hai thành viên trong nhóm phát triển.
        \item \textbf{KPI và Tiêu chí hoàn thành:} Hoàn thiện tập dữ liệu mẫu chuẩn; xây dựng script tự động chạy đo lường Precision, Recall và F1-score của mô hình phân loại giao dịch.
    \end{itemize}

    \item \textbf{Đánh giá hiệu năng hạ tầng bằng kiểm thử tải:}
    \begin{itemize}
        \item \textbf{Độ ưu tiên:} Trung bình.
        \item \textbf{Nguồn lực dự kiến:} 01 DevOps/QA Engineer, máy chủ chạy thử tải (sử dụng k6 hoặc Artillery).
        \item \textbf{Sự phụ thuộc:} Không có.
        \item \textbf{Rủi ro và Giải pháp:} Quá trình kiểm thử tải lớn có thể làm nghẽn hoặc sập hệ thống dịch vụ chính. Giải pháp là triển khai trên một môi trường thử nghiệm phụ (sandbox) cô lập với môi trường chính.
        \item \textbf{KPI và Tiêu chí hoàn thành:} Đo lường được ngưỡng chịu tải tối đa của hệ thống thông báo đẩy (FCM); tỷ lệ lỗi dưới 1\% ở mức tải thiết kế thông thường.
    \end{itemize}
\end{enumerate}

\subsection{Next: Mở rộng khả năng tự động hóa và tích hợp}

\begin{enumerate}
    \item \textbf{Hoàn thiện tính năng quét ảnh hóa đơn (OCR Receipt):}
    \begin{itemize}
        \item \textbf{Độ ưu tiên:} Trung bình.
        \item \textbf{Nguồn lực dự kiến:} 01 Mobile Developer, API nhận dạng ký tự quang học (như Google Cloud Vision OCR).
        \item \textbf{Sự phụ thuộc:} Hoàn thành Database Transaction cho luồng tài chính.
        \item \textbf{Rủi ro và Giải pháp:} Hóa đơn thực tế bị nhăn, mờ hoặc thiếu sáng làm giảm độ chính xác khi đọc chữ. Giải pháp là thiết kế giao diện hướng dẫn người dùng chụp ảnh trực diện, đủ sáng và cắt bớt phần thừa của hóa đơn.
        \item \textbf{KPI và Tiêu chí hoàn thành:} Trích xuất đúng các thông tin cơ bản (ngày giao dịch, số tiền) đạt độ chính xác trên 80\% đối với các mẫu hóa đơn rõ nét.
    \end{itemize}

    \item \textbf{Mở rộng kết nối API ngân hàng:}
    \begin{itemize}
        \item \textbf{Độ ưu tiên:} Cao.
        \item \textbf{Nguồn lực dự kiến:} 01 Backend Developer, cổng dịch vụ Open Banking trung gian (như Casso hoặc VietQR).
        \item \textbf{Sự phụ thuộc:} Đánh giá đầy đủ các yêu cầu pháp lý và bảo mật thông tin (chi tiết tại Mục~\ref{subsec:key-items-analysis}).
        \item \textbf{Rủi ro và Giải pháp:} API của ngân hàng gặp sự cố gián đoạn hoặc bảo trì. Giải pháp là sử dụng hàng đợi BullMQ để lưu vết các yêu cầu đồng bộ và tự động thực hiện lại khi kết nối được phục hồi.
        \item \textbf{KPI và Tiêu chí hoàn thành:} Đồng bộ tự động các giao dịch từ tài khoản ngân hàng liên kết thử nghiệm với số dư trong ứng dụng di động.
    \end{itemize}

    \item \textbf{Bổ sung chức năng phát hiện giao dịch bất thường:}
    \begin{itemize}
        \item \textbf{Độ ưu tiên:} Thấp.
        \item \textbf{Nguồn lực dự kiến:} 01 Backend/AI Developer, máy chủ cơ sở dữ liệu.
        \item \textbf{Sự phụ thuộc:} Hoàn thành bộ benchmark AI và tích lũy dữ liệu chi tiêu thực tế tối thiểu của người dùng thử nghiệm.
        \item \textbf{Rủi ro và Giải pháp:} Thuật toán gửi quá nhiều cảnh báo sai gây phiền hà cho người dùng. Giải pháp là cho phép người dùng phản hồi xác nhận giao dịch bình thường để tinh chỉnh thuật toán cá nhân hóa.
        \item \textbf{KPI và Tiêu chí hoàn thành:} Mô hình phát hiện được các khoản chi tiêu vượt ngưỡng bất thường và gửi cảnh báo thành công qua thông báo đẩy.
    \end{itemize}
\end{enumerate}

\subsection{Later: Nâng cao năng lực AI}

\begin{enumerate}
    \item \textbf{Nâng cấp cơ chế gợi ý học bổng bằng Semantic Search:}
    \begin{itemize}
        \item \textbf{Độ ưu tiên:} Trung bình.
        \item \textbf{Nguồn lực dự kiến:} 01 AI Engineer, cơ sở dữ liệu hỗ trợ lưu trữ vector (pgvector).
        \item \textbf{Sự phụ thuộc:} Chuẩn hóa dữ liệu học bổng đầu vào ở giai đoạn Now.
        \item \textbf{Rủi ro và Giải pháp:} Tìm kiếm theo ngữ nghĩa (Semantic Search) yêu cầu nhiều tài nguyên tính toán hơn tìm kiếm từ khóa. Giải pháp là lập chỉ mục vector để giảm thời gian tìm kiếm.
        \item \textbf{KPI và Tiêu chí hoàn thành:} Nhận diện được các trường hợp đồng nghĩa trong hồ sơ (ví dụ: "Kỹ thuật phần mềm" tương đương với "Công nghệ thông tin" trong một số ngữ cảnh học bổng) với thời gian phản hồi dưới 1 giây.
    \end{itemize}

    \item \textbf{Tối ưu hóa mô hình AI nội bộ (Local LLM):}
    \begin{itemize}
        \item \textbf{Độ ưu tiên:} Trung bình.
        \item \textbf{Nguồn lực dự kiến:} 01 AI Infrastructure Engineer, máy chủ vật lý hoặc dịch vụ đám mây hỗ trợ GPU chuyên dụng.
        \item \textbf{Sự phụ thuộc:} Hoàn thành Semantic Search và kiểm thử tải hiệu năng ở các giai đoạn trước.
        \item \textbf{Rủi ro và Giải pháp:} Mô hình chạy nội bộ phản hồi chậm khi có nhiều yêu cầu đồng thời. Giải pháp là lượng tử hóa mô hình (Quantization 4-bit) để tiết kiệm bộ nhớ và sử dụng vLLM để cải thiện tốc độ sinh token.
        \item \textbf{KPI và Tiêu chí hoàn thành:} Triển khai thành công mô hình ngôn ngữ lớn trên máy chủ nội bộ của trường đại học phục vụ hội thoại tư vấn cơ bản; thời gian phản hồi trung bình dưới 2 giây.
    \end{itemize}
\end{enumerate}

\subsection{Phân tích tính khả thi cho các hạng mục trọng điểm}
\label{subsec:key-items-analysis}

Để chuẩn bị triển khai các mục tiêu lớn trong lộ trình phát triển bao gồm liên kết tài khoản ngân hàng thực tế và tối ưu hóa mô hình trí tuệ nhân tạo chạy nội bộ (Local LLM), nhóm phát triển thực hiện phân tích chi tiết về khía cạnh hạ tầng, pháp lý, chi phí và phạm vi MVP cho từng hạng mục như sau:

\textbf{1. Liên kết tài khoản ngân hàng thực tế}
\begin{itemize}
    \item \textbf{Hạ tầng kỹ thuật:} Hệ thống sẽ kết nối với cổng dịch vụ Open Banking trung gian tại Việt Nam qua giao thức HTTPS bảo mật bằng mã hóa SSL/TLS và chữ ký số. Cổng trung gian này cung cấp webhook để gửi thông tin biến động số dư thời gian thực. Hệ thống sử dụng một dịch vụ API Gateway nội bộ để tiếp nhận, xác thực chữ ký và đẩy vào hàng đợi BullMQ để xử lý bất đồng bộ, tránh nghẽn luồng dữ liệu.
    \item \textbf{Khía cạnh pháp lý:} Việc lưu trữ và xử lý thông tin giao dịch tài chính phải tuân thủ nghiêm ngặt Nghị định số 13/2023/NĐ-CP về bảo vệ dữ liệu cá nhân tại Việt Nam. Hệ thống bắt buộc phải có sự đồng ý rõ ràng của người dùng (Opt-in) trước khi liên kết, mã hóa thông tin tài khoản ngân hàng trong cơ sở dữ liệu và chỉ truy cập lịch sử giao dịch ở chế độ đọc (Read-only), tuyệt đối không lưu trữ mật khẩu hoặc mã PIN ngân hàng của người dùng.
    \item \textbf{Chi phí dự kiến:} Tận dụng các gói dịch vụ Open Banking dành cho nhà phát triển có mức phí duy trì thấp (khoảng 100.000 -- 300.000 VNĐ/tháng). Chi phí này được chi trả từ nguồn quỹ nghiên cứu khoa học của nhà trường hoặc mô hình xã hội hóa dịch vụ.
    \item \textbf{Phạm vi phiên bản tối thiểu (MVP):} Ở giai đoạn đầu, hệ thống chỉ hỗ trợ liên kết với một hoặc hai ngân hàng phổ biến có cung cấp API nhận biến động số dư miễn phí hoặc giá rẻ (ví dụ: qua cổng dịch vụ tích hợp VietQR). Phiên bản MVP chỉ dừng lại ở việc đọc lịch sử biến động số dư để đối soát tự động với các khoản chi tiêu, chưa hỗ trợ thực hiện các lệnh chuyển tiền hoặc thanh toán trực tiếp từ ứng dụng.
\end{itemize}

\textbf{2. Triển khai mô hình AI nội bộ (Local LLM)}
\begin{itemize}
    \item \textbf{Hạ tầng kỹ thuật:} Để phục vụ mô hình ngôn ngữ lớn chạy nội bộ ổn định với thời gian phản hồi nhanh, hệ thống yêu cầu một máy chủ vật lý được trang bị tối thiểu 01 GPU chuyên dụng (ví dụ: NVIDIA RTX 4090 24GB VRAM hoặc NVIDIA L4 24GB VRAM). Mô hình sử dụng các framework phục vụ hiệu năng cao như vLLM hoặc Ollama kết hợp với kỹ thuật lượng tử hóa 4-bit (4-bit quantization) để giảm dung lượng mô hình, giúp chạy vừa trong bộ nhớ VRAM của một GPU đơn lẻ.
    \item \textbf{Khía cạnh pháp lý:} Do mô hình chạy hoàn toàn trên máy chủ đặt tại phòng máy của nhà trường, toàn bộ dữ liệu hội thoại của sinh viên không bị gửi lên môi trường đám mây của bên thứ ba. Điều này giúp đảm bảo an toàn thông tin theo các quy định nội bộ của nhà trường và hạn chế rò rỉ thông tin cá nhân của sinh viên ra ngoài.
    \item \textbf{Chi phí dự kiến:} Chi phí mua sắm thiết bị ban đầu (khoảng 50 -- 70 triệu VNĐ cho một máy chủ GPU tầm trung) hoặc thuê máy chủ ảo GPU trên đám mây (khoảng 1.500.000 -- 3.000.000 VNĐ/tháng). Chi phí vận hành dài hạn thấp do không phải trả tiền theo số lượng ký tự (token) sử dụng như các dịch vụ Cloud API thương mại.
    \item \textbf{Phạm vi phiên bản tối thiểu (MVP):} Phiên bản MVP sử dụng các mô hình ngôn ngữ lớn nguồn mở có kích thước nhỏ và phổ biến (ví dụ: Qwen-2.5-7B-Instruct hoặc Llama-3-8B-Instruct đã được lượng tử hóa). Trợ lý ảo trong phiên bản MVP chỉ tập trung xử lý các câu hỏi thường gặp về quy chế học thuật của nhà trường và phân loại giao dịch cơ bản, các yêu cầu phân tích tài chính phức tạp hơn vẫn được định tuyến (route) tạm thời qua Cloud API (như Gemini) để đảm bảo chất lượng câu trả lời.
\end{itemize}
